\documentclass[10pt,a4paper]{amsart}
\usepackage[margin=19mm]{geometry}
\usepackage{amsmath,amssymb,mathtools}
\usepackage[scr=rsfs,cal=euler]{mathalfa}
\usepackage{xfrac}
\usepackage[unicode,hidelinks,bookmarks=false]{hyperref}
\newcommand{\ie}{i.e.\ }
\newcommand{\cf}{cf.\ }
\newcommand{\resp}{respectively\ }
\newcommand{\Subsection}{\subsection}
\providecommand{\nobreakdash}{\nobreak\hbox{-}}
\setlength{\emergencystretch}{2em}
\multlinegap0pt
\usepackage{xcolor}
\usepackage{amssymb}
\usepackage{enumerate}
\usepackage{bm}
\usepackage{float}
\usepackage{csquotes}
\usepackage{tikz}
\newtheorem{lemma}{Lemma}
\title{On the cycle structure of the symmetric tensor power of permutations}
\author{Sebastian Caballero, Diego Villamizar}
\date{Source: arXiv:2605.27214v1 (CC BY 4.0)}
\begin{document}
\maketitle

\noindent\emph{Source and licence.} On the cycle structure of the symmetric tensor power of permutations. Version arXiv:2605.27214v1 (CC BY 4.0), distributed by arXiv under the Creative Commons Attribution 4.0 International licence, http://creativecommons.org/licenses/by/4.0/, as recorded in the arXiv licence metadata for this exact version and shown on the abstract page https://arxiv.org/abs/2605.27214v1. Full licence notice: \texttt{licenses/arxiv-2605.27214-CC-BY-4.0.txt}.\par\smallskip

\noindent\textit{Editorial scope.} Counted unit: the article's lemma lemma:eigenvalues (source0.tex lines 230--233). The article's complete original proof of it is delivered with it (source0.tex lines 235--245, one \texttt{proof} environment of 11 lines). No mathematics is written, repaired or re-derived: the statement and the proof are the article's own lines, in the article's own notation, and no retained line is substituted or re-flowed.  No further source-local context is needed: every cross-reference of every retained line resolves inside the two delivered spans.  Nothing was omitted from any retained span, so the package carries no omission record.  No retained line contains a citation, so the wrapper carries no bibliography and no bibliography entry.  No retained line contains a figure environment, an image-inclusion command or drawing code, so no image is delivered and the package declares no asset dependency.  Editorial formatting only, in addition: an AMS \texttt{amsart} preamble that restates the article's own declarations and macros used by the retained lines, namely the environments \texttt{lemma} on separate counters, together with the standard packages those lines need.  The retained environments print in this document as Lemma 1 in order of appearance, whereas the article prints the same items with the numbering of its own full document.  The complete native source of the article as distributed in the arXiv e-print for this exact version is bundled unchanged as \texttt{sources/201474/source0.tex}.
    \begin{lemma}
    \label{lemma:eigenvalues}
        Let $\sigma \in \mathfrak{S}_n$ and $\ell \in [n]$. A cycle of length $\ell$ in the decomposition of $\sigma$ contributes once to the multiplicity of the eigenvalue $e^{\frac{2\pi i k}{\ell}}$ for $0 \le k < \ell-1$.
    \end{lemma}

    \begin{proof}
        Let $\sigma = (a_1\,a_2\, \dots \, a_\ell)$ be any cycle of length $\ell$, and let $k$ be a non-negative integer $0 \le k < \ell-1$. We explicitly construct an eigenvector for $\lambda = e^{\frac{2 \pi i k}{\ell}}$. Let $\mathbf{v} = (v_1, v_2, \dots, v_n)$ be defined as $\mathbf{v} = \sum _{i = 1}^{\ell}\lambda ^ {i-1}e_{a_{i}}$, where $e_i = (0,\cdots ,1,\cdots ,0)$ denotes the $i$-th element of the standard basis. $\textbf{v}$ contains powers of $\lambda$ exactly at the positions of the elements that belong to the cycle.
        %\begin{align*}
        %    v_j &= 0 \text{ for $j$ not in the cycle}\\ 
        %    v_{a_1} &= 1\\
        %    v_{a_{i+1}} &= \lambda v_{a_i} \, , \text{ for }1<i<\ell %.
        %\end{align*}
        This construction implies that $\lambda$ is an eigenvalue for $A_\sigma$ with eigenvector $\mathbf{v}$ noticing that for $j$ not in the cycle, $(A_\sigma \mathbf{v})_j = 0 = (\lambda \mathbf{v})_j$; for $i \in [\ell-1]$, we have $(A\mathbf{v})_{a_{i}} = v_{a_{i+1}} = (\lambda \mathbf{v})_{a_i}$; and for the case $i = \ell$, we have $(A\mathbf{v})_{a_\ell} = \mathbf{v}_1 = 1 = (\lambda \mathbf{v})_{a_{\ell}}$ 
        therefore, $A\mathbf{v} = \lambda \mathbf{v}$.
        Each cycle of length $\ell$ contributes exactly $\ell$ eigenvalues. All eigenvectors corresponding to different cycles are linearly independent, even if their eigenvalues coincide, so the set of all eigenvectors has size $\sum _{\ell = 1}^n \ell \cdot C_{\ell}(\sigma) = n$. This set is a basis for $\mathbb{C}^n$, ensuring that each cycle contributes exactly once for the multiplicity of each of its eigenvalues.
    \end{proof}

\end{document}
